% ============================================================================
\clearpage
\section*{v4.10.0-r03 Source-Only: Fermions and Bosons at the Big Bang, and the Vacuum Term under G2}
\addcontentsline{toc}{section}{v4.10.0-r03 Source-Only: Fermions and Bosons at the Big Bang, and the Vacuum Term under G2}
\label{sec:v4100-r03-bang-g2}
% ============================================================================

\paragraph{Purpose.}
Layer v4.10.0-r02 closed the asymptotics of the bosonic admixture at the big bang (\S\ref{sec:v4100-r02-admixture-asymptotics}).
This layer records what follows for the other fields and for the vacuum term under the domain choice G2 (tachyon reading before,
matter reading after the bang; r85).  It contains one pre-registered test (H94), an author decision on the duration \(\Delta\) of the
transition, an exact symmetry of the two senses of rotation, and analytic explorations of four routes to \(\Lambda\) under G2.  Two of my own
intermediate claims are corrected here.  The notes are
\texttt{notes/\allowbreak research/\allowbreak v4\_10\_0\_r03/}.  Earlier layers are not rewritten.

\noindent\textbf{Source-only status.}
This v4.10.0-r03 layer is source-only.  It does not build a PDF, does not assign a new DOI,
and does not authorize a public release.  The v4.10.0 Version DOI
\href{https://doi.org/10.5281/zenodo.23991357}{10.5281/zenodo.23991357} and the Concept DOI
\href{https://doi.org/10.5281/zenodo.19210728}{10.5281/zenodo.19210728} are unchanged by it.

\subsection*{v4.10.0-r03 results}
\label{subsec:v4100-r03-results}

\paragraph{Theoremlet v4.10.0-r03.1 (fermions pass the bang without admixture).}
Exploration (analytic, no criterion).  Under G2 and the second reading the symmetric Dirac or Weyl action is weighted by \(c(\tau)\), going from
\(-b\) to \(a\).  The field equation \(c(i\gamma^\mu\partial_\mu-m)\psi+\tfrac i2c'\gamma^0\psi=0\) is solved exactly by \(\psi=c^{-1/2}\chi\), with \(\chi\) free, for
every profile (sympy, residual \(0\)).
\begin{enumerate}[leftmargin=2em]
\item[(a)] Positive and negative frequencies do not mix: no fermions are created at the bang, for any \(k\) and any profile.  A first-order equation lets
the weight be scaled away; the second-order equations of scalars and gauge fields do not.
\item[(b)] Transfer factor \(\psi_{\rm after}/\psi_{\rm before}=\pm i\sqrt{b/a}=\pm i\beta_{\rm sym}\) (continuous branch along the path; the two senses differ in phase
only).  The map is unitary in the weighted norm, and the frequency is kept: a particle of the T reading becomes an antiparticle-like state of the M
reading, exactly and for all \(k\).
\item[(c)] Correction (made in the notes the same day): mass terms keep their sign under the image (r85, H82), so before the bang the mass is \(-m\)
relative to the kinetic term.  On FRW the mass term is \(m\,a(\tau)\) with \(a\to0\) at the bang, so (a) and (b) hold for \(k\gg m\,a\) at the transition.
\end{enumerate}

\paragraph{Theoremlet v4.10.0-r03.2 (H94: the weighted norm at the bang).}
Pre-registered (commit \texttt{99f1cce}, partly not blind).  \(N=a(|C|^2-|D|^2)/b\) for the mode of positive norm before the bang (tanh profile,
\(\Delta=2\)).  All four parts confirmed (one run, Wronskian constant to \(10^{-11}\)).
\begin{enumerate}[leftmargin=2em]
\item[(a)] Bosons: \(N\to+1\) for long waves (\(1.024\), \(0.974\) at \(k\Delta=0.01\)) and \(N\to-1\) for short waves (\(-0.969\), \(-1.032\) at \(20\)).
\item[(b)] Exactly one sign change per sense of rotation, at \(k^*\Delta=0.7361\) and \(0.2569\).
\item[(c)] In between the norm is not conserved (\(0.1<|N|<0.9\) at \(10\) and \(5\) grid points).
\item[(d)] Fermions: \(N_F=-1\) on all points, to \(6\times10^{-12}\).
\end{enumerate}
\emph{Noticed afterwards.}  The bosonic norm is also amplified (up to \(1.18\) and \(-1.63\)).  The mean of the two senses does not restore it.  The
statement ``bosons with matter weights are created, \(n=0.404\) per mode'' holds only in the picture of the sharp real jump.

\paragraph{Decision v4.10.0-r03.3 (author, 6 November 2026: \(\Delta\lesssim\ell_P\)).}
The duration of the transition is taken at or below the Planck length.  All modes below the Planck scale then have \(k\Delta\ll1\): bosons with matter
weights pass norm-preserving as a Bogoliubov transformation with \(n=(1-\beta_{\rm sym}^2)^2/(4\beta_{\rm sym}^2)=0.404\) per mode, and fermions are mirrored
by CPT without creation.  The non-unitary range \(k\sim1/\Delta\) lies at or beyond the Planck scale.  The total energy of the created bosons is set by
that range and stays open.

\paragraph{Theoremlet v4.10.0-r03.4 (the two senses of rotation are exactly conjugate).}
Exploration (derived analytically, checked numerically).  For the transfer matrix \(M(L)=\begin{pmatrix}a&b\\c&d\end{pmatrix}\) in the basis
\((e^{ikt},e^{-ikt})\) one has \(\det M=e^{-L}\) (weighted Wronskian).  The duality \(v=wy'\) gives \(M(-L)=e^L JMJ\) with \(J=\operatorname{diag}(1,-1)\).  Reflection
with conjugation, using \(s(-t)=1-s(t)\), gives \(M(-\bar L)=\overline{M(L)}^{-1}\).  Together
\[
M(\bar L)=\begin{pmatrix}\bar d&\bar b\\ \bar c&\bar a\end{pmatrix}.
\]
Hence \(C(\bar L)=\overline{C(L)}\), so \(|C_+|=|C_-|\) exactly, and \(D_+\overline{D_-}-b_+C_+=e^{-L}\), the product rule of r85.6 with its correction term.  For
the local problem of r02.2, \(M=\begin{pmatrix}1&0\\r&1\end{pmatrix}\) exactly, so \(r(\bar\lambda)=\overline{r(\lambda)}\).  Numerically all identities hold to
\(\le3\times10^{-10}\) at \(k\Delta=0.01\)--\(5\).

\paragraph{Theoremlet v4.10.0-r03.5 (no \(\Lambda\) from the tachyon sector after the bang under G2).}
Exploration (analytic, no criterion).  Four routes were examined.
\begin{enumerate}[leftmargin=2em]
\item[(a)] Mixture (deferred idea I3): its \(\Lambda\) side came from the tension of the T reading.  Under G2 only the M reading exists after the bang, so
only the dust part remains.
\item[(b)] Tachyonic condensation (I4): after the bang there are no tachyon modes.  Before it the image of a scalar sees \(-V\) (\(\mathcal L_T=-\tfrac12X-V\),
r84); a potential stable in M runs away there (\(\ddot\phi=m^2\phi+\lambda\phi^3\)) instead of condensing.  Under G1 one has \(m_{\rm eff}^2=\gamma m^2>0\).
\item[(c)] Ghost condensate (r83.4): under G1, \(aP_M+bP_T=(a-b)X/2+(a+b)cX^2\) is healthy and does not condense.  Under G2 the condensate exists only
before the bang; after it \(P_M\) has \(w\in[\tfrac13,1]\), never \(-1\); at the bang the condensate has \(\pi=0\), so the M field is at rest.
\item[(d)] Conclusion: under G2 a vacuum term after the bang can only be \(aV_{\min}\) of the M reading, with a free value.  The structure does not fix
\(\Lambda\) under G2.
\end{enumerate}

\paragraph{Theoremlet v4.10.0-r03.6 (boundary condition and spatial reflection; correction).}
Exploration (analytic).  If the T field before the bang is identified pointwise with the image of the M field after it,
\(\phi_T(x)=\phi_M(\Lambda(i\pi)x)\), the junction yields \(\phi'(0)=0\) for scalars and \(\psi(0)=0\) for fermions (\(\beta_{\rm sym}\neq1\)).  The same identification
needs a direction \(n\) and a spatial origin (from \(x\mapsto\Lambda x+2x_0\)) that nothing fixes; classically it would correlate the modes \(k\) and \(P_nk\).
Quantum mechanically it is inconsistent for \(a\neq b\): with antiunitary time reversal the junction forces \(u'(0)=0\), and the Wronskian
normalization \(uu'^*-u^*u'=i\) then fails at \(\tau=0\).  \emph{Correction:} the conditions \(\phi'(0)=0\), \(\psi(0)=0\) stated the same day in the notes do
not follow from the structure; a frozen scalar at the bang is a possible choice of state, not a consequence.  Under G2 the second reading is a
statement about the action, \(-b\mathcal L_M\) before and \(a\mathcal L_M\) after, with independent data before the bang, as in H85--H94; at that level
neither a direction nor an origin appears (r84, H79).

\paragraph{Proof and measurement sources.}
Scripts \texttt{k209\_\allowbreak fermionen\_\allowbreak urknall.py}, \texttt{k210\_\allowbreak norm\_\allowbreak uebergang.py},
\texttt{k211\_\allowbreak symmetrie\_\allowbreak drehsinne.py} with their outputs, and the explorations in \texttt{notes/\allowbreak research/\allowbreak v4\_10\_0\_r03/};
pre-registration and result H94 in \texttt{notes/\allowbreak research/\allowbreak predictions/}.  \(\square\)

\paragraph{Consequence.}
Under G2 the big bang separates bosons from fermions.  Fermions are mirrored by CPT exactly and unitarily.  Bosons with matter weights are
Bogoliubov-transformed for \(k\Delta\ll1\) and CPT-mirrored for \(k\Delta\gg1\), and the transition is non-unitary in between, which the decision
\(\Delta\lesssim\ell_P\) moves to the Planck scale.  The two senses of rotation are exactly conjugate.  None of the tachyon-sector mechanisms yields a
vacuum term after the bang; \(\Lambda\) remains a free \(aV_{\min}\).  OP5 is not closed.

\subsection*{v4.10.0-r03 status in the three frames}
\label{subsec:v4100-r03-frames}

\begin{itemize}[leftmargin=2em]
\item r03.1--r03.4 concern the M--T transition under G2, with the weights \(a\), \(b\) of frame M: frame T before and frame M after the bang.  No frame-free
statement is made.
\item Photon frame P: with \(a=b\) there is no admixture, the fermion factor has modulus \(1\), and the pointwise identification of r03.6 is not obstructed.
The norm then stays \(+1\) for bosons, as in r85.5.
\item r03.5 states that frame T contributes no vacuum term in frame M after the bang under G2.  Under G1 both readings coexist, and the reduced theory is
the M theory with renormalized weights (r84).
\end{itemize}

\subsection*{v4.10.0-r03 register entries}
\label{subsec:v4100-r03-register}

\begin{center}
\small
\begin{tabular}{@{}>{\raggedright\arraybackslash}p{0.34\textwidth}>{\raggedright\arraybackslash}p{0.58\textwidth}@{}}
\toprule
statement & status after v4.10.0-r03\\
\midrule
fermions at the bang (G2) & no admixture, CPT exact, factor \(\pm i\beta_{\rm sym}\) (exploration)\\
bosonic norm at the bang & \(+1\to-1\) with one zero; non-unitary near \(k\sim1/\Delta\) (H94)\\
duration \(\Delta\) & \(\lesssim\ell_P\) (author decision)\\
senses of rotation & exactly conjugate, \(|C_+|=|C_-|\) (exploration)\\
\(\Lambda\) from I3, I4, ghost condensate under G2 & none after the bang; \(\Lambda=aV_{\min}\) free\\
pointwise identification of T and M & inconsistent for \(a\neq b\); boundary conditions withdrawn\\
\bottomrule
\end{tabular}
\end{center}

\subsection*{v4.10.0-r03 ledger and guard update}
\label{subsec:v4100-r03-ledger-guard}

\[
R_{\rm chk}^{v4100\text{-}r03}=(1,1,0)
\]
for (fermions and the bosonic norm at the bang settled; symmetry exact; \(\Lambda\) under G2 without a mechanism).

The layer spends no guard:
\[
R_{\rm glob}^{(v4100\text{-}r03,+)}=(0,1,1,1,1),
\qquad
\Gamma_{\rm guard}^{v4100\text{-}r03}=(1,1,1,1,1),
\qquad
\beta_4=1/256\ \text{unspent}.
\]
It does not imply exact finite \(C5\) zero closure, exact all-readout terminality,
\(C6/T2\), a global photon-clock/RH theorem, or global all-readout closure.
